\documentclass[11pt,oneside]{article}

\usepackage[a4paper,margin=27mm,headheight=15pt]{geometry}
\usepackage[T1]{fontenc}
\usepackage[utf8]{inputenc}
\IfFileExists{libertinus.sty}{\usepackage{libertinus}}{\usepackage{lmodern}}
\usepackage{microtype}
\usepackage{amsmath,amssymb,amsthm,mathtools,mathrsfs}
\usepackage{bm}
\usepackage{booktabs,longtable,array,tabularx}
\usepackage{graphicx}
\usepackage{pgfplots}
\usepackage{xcolor}
\usepackage{enumitem}
\usepackage{fancyhdr}
\usepackage{titlesec}
\usepackage{listings}
\usepackage{xurl}
\usepackage{hyperref}
\usepackage[nameinlink,noabbrev]{cleveref}

\definecolor{linkblue}{RGB}{25,75,135}
\definecolor{shadegray}{RGB}{246,247,249}
\definecolor{rulegray}{RGB}{195,201,208}
\pgfplotsset{compat=1.18}
\hypersetup{
  colorlinks=true,
  linkcolor=linkblue,
  citecolor=linkblue,
  urlcolor=linkblue,
  pdftitle={Further Exact Formulae for the Thue--Morse Sequence},
  pdfsubject={Digit, valuation, Boolean, generating-function, finite-difference, and Fourier formulae},
  pdfkeywords={Thue-Morse, binary digit sum, Pascal triangle, Lucas theorem, Kummer theorem, Prouhet, Walsh transform}
}

\pagestyle{fancy}
\fancyhf{}
\fancyhead[L]{\small Further Thue--Morse formulae}
\fancyhead[R]{\small\nouppercase{\leftmark}}
\fancyfoot[C]{\thepage}
\renewcommand{\headrulewidth}{0.4pt}

\titleformat{\section}{\Large\bfseries}{\thesection}{0.7em}{}
\titleformat{\subsection}{\large\bfseries}{\thesubsection}{0.7em}{}
\titleformat{\subsubsection}{\normalsize\bfseries}{\thesubsubsection}{0.7em}{}
\setcounter{secnumdepth}{3}
\setcounter{tocdepth}{2}
\setlist[itemize]{leftmargin=2em,itemsep=0.25em,topsep=0.4em}
\setlist[enumerate]{leftmargin=2.3em,itemsep=0.25em,topsep=0.4em}
\setlength{\emergencystretch}{2em}

\newtheorem{theorem}{Theorem}[section]
\newtheorem{proposition}[theorem]{Proposition}
\newtheorem{lemma}[theorem]{Lemma}
\newtheorem{corollary}[theorem]{Corollary}
\newtheorem{conjecture}[theorem]{Conjecture}
\theoremstyle{definition}
\newtheorem{definition}[theorem]{Definition}
\newtheorem{algorithm}[theorem]{Algorithm}
\newtheorem{example}[theorem]{Example}
\theoremstyle{remark}
\newtheorem{remark}[theorem]{Remark}
\newtheorem{warning}[theorem]{Warning}

\newcommand{\R}{\mathbb R}
\newcommand{\C}{\mathbb C}
\newcommand{\Q}{\mathbb Q}
\newcommand{\N}{\mathbb N}
\newcommand{\Z}{\mathbb Z}
\newcommand{\wt}{\operatorname{wt}_2}
\newcommand{\vtwo}{\nu_2}
\newcommand{\sgn}{\operatorname{sgn}}
\newcommand{\e}{\mathrm e}
\newcommand{\ii}{\mathrm i}
\newcommand{\dd}{\,\mathrm d}
\newcommand{\ind}{\mathbf 1}
\newcommand{\xor}{\mathbin{\oplus}}
\newcommand{\preceqbits}{\mathrel{\preccurlyeq_2}}
\newcommand{\bitand}{\mathbin{\&}}
\newcommand{\Bell}{\mathcal B}
\newcommand{\bigO}{\mathcal O}
\newcommand{\smallo}{o}
\newcommand{\repo}{\nolinkurl{Analysis/FabiusFunction}}
\newcommand{\lean}[1]{%
  \ifmmode\text{\nolinkurl{#1}}\else\nolinkurl{#1}\fi}

\newenvironment{proofidea}{\par\smallskip\noindent\textit{Proof architecture.}\ }{\hfill$\square$\par\smallskip}
\newenvironment{boxedremark}{%
  \par\smallskip\noindent\begin{center}\begin{minipage}{0.94\textwidth}%
  \color{black}\hrule\vspace{0.6em}\small
}{%
  \vspace{0.6em}\hrule\end{minipage}\end{center}\smallskip
}

\lstdefinestyle{pseudo}{
  basicstyle=\ttfamily\small,
  backgroundcolor=\color{shadegray},
  frame=single,
  rulecolor=\color{rulegray},
  columns=fullflexible,
  keepspaces=true,
  showstringspaces=false,
  xleftmargin=0.7em,
  xrightmargin=0.7em,
  aboveskip=0.8em,
  belowskip=0.8em
}

\title{\bfseries Further Exact Formulae for the Thue--Morse Sequence\\[0.35em]
\large Digit, valuation, Boolean, finite-difference, Fourier, and analytic representations}
\author{Companion report to the Fabius--Rvachev formal development\\
\small Vladimir Reshetnikov's \texttt{ProveIt} repository}
\date{25 August 2026\\
\small source comparison at commit \nolinkurl{eaea1b9afe2b7ff0b7dd880bd1710e303dec6d80}}

\begin{document}
\maketitle

\begin{abstract}
Section~11.8 of the companion Fabius--Rvachev article recovers the zero--one
Thue--Morse bit from the parity pattern of one row of Pascal's triangle by an
exact logarithm.  The present report searches systematically for structurally
different exact formulae.  The most economical arithmetic answer is
\[
  \tau(n)\equiv \vtwo\binom{2n}{n}\pmod2,
  \qquad
  \varepsilon_n=(-1)^{\vtwo\binom{2n}{n}},
\]
so a single central binomial coefficient replaces the whole Pascal row.  Other
formulae express the same sign through dyadic floor sums, the parity of
$\vtwo(n!)$, binary carry counts, XOR characters, the Möbius function of the
Boolean submask lattice, and a finite trigonometric sign product.

The report then develops several less immediate representations.  Lucas's
theorem produces a Pascal-parity Möbius inversion and a finite ``binomial only''
formula.  The lacunary Euler product gives a logarithmic-derivative recurrence,
an explicit partition sum, an exponential Bell-polynomial formula, and a lower
Hessenberg determinant for each term.  A master mixed finite-difference identity
contains Prouhet cancellation, the first surviving moment, exponential products,
and positive reciprocal-power sums as special cases.  Finally, finite discrete
Fourier inversion gives a cyclotomic formula for an individual bit, while the
Mellin transform gives an entire continuation of the associated shifted
Dirichlet series.  A short final section records the Woods--Robbins product and
a proposed Lean formalization map.

Every mathematical assertion below is proved in the text.  Except where an
existing declaration is explicitly named, no claim is made that the new
formulae are already formalized in the repository.
\end{abstract}

\tableofcontents
\newpage

\section{Scope, conventions, and the formula being complemented}
\label{sec:scope}

Write the binary expansion of $n\in\N$ as
\begin{equation}
  n=\sum_{j\ge0} b_j(n)2^j,
  \qquad b_j(n)\in\{0,1\},                                  \label{eq:binary-expansion}
\end{equation}
with only finitely many nonzero digits.  Its binary weight and the two standard
normalizations of the Thue--Morse sequence are
\begin{equation}
  \wt(n)=\sum_{j\ge0}b_j(n),
  \qquad
  \tau(n)=\wt(n)\bmod2\in\{0,1\},
  \qquad
  \varepsilon_n=(-1)^{\wt(n)}=1-2\tau(n).                  \label{eq:basic-definitions}
\end{equation}
The sign normalization is algebraically cleaner; a formula for one normalization
immediately gives a formula for the other through
\begin{equation}
  \tau(n)=\frac{1-\varepsilon_n}{2},
  \qquad
  \varepsilon_n=1-2\tau(n).                                \label{eq:normalization-bridge}
\end{equation}

The source article's section ``A binomial--logarithm formula for the
Thue--Morse bit'' sets
\begin{equation}
  \mathcal X(n)=\sum_{k=0}^{n}(-1)^{\binom nk}              \label{eq:source-X}
\end{equation}
and proves
\begin{equation}
  n+1-\mathcal X(n)=2^{\wt(n)+1}.                           \label{eq:source-exact-power}
\end{equation}
Consequently
\begin{equation}
  \tau(n)=
  \frac{1+(-1)^{\log_2\!\left(n+1-
       \sum_{k=0}^{n}(-1)^{\binom nk}\right)}}{2}.          \label{eq:source-binomial-log}
\end{equation}
The exact-power theorem is essential: the displayed logarithm is not being used
as a numerical approximation.  Formula~\eqref{eq:source-binomial-log} hides the
binary representation inside the odd entries of a single Pascal row.  Our goal
is not merely to rewrite it, but to locate other structures in which the same
parity character appears naturally.

We use the following notation.
\begin{itemize}
\item $\vtwo(m)$ is the exponent of $2$ in a positive integer $m$, with
      $\vtwo(1)=0$.
\item $a\xor b$ is bitwise exclusive-or.
\item $k\preceqbits n$ means that every $1$-bit of $k$ is also a $1$-bit of $n$;
      equivalently, $k\bitand\neg n=0$.  We call $k$ a \emph{submask} of $n$.
\item $[P]$ or $\ind_P$ denotes the indicator of a proposition $P$.
\item Empty products are $1$ and empty sums are $0$.
\end{itemize}

\begin{table}[ht]
\centering
\small
\begin{tabular}{@{}rrrrrrrr@{}}
\toprule
$n$ & binary & $\wt(n)$ & $\tau(n)$ & $\varepsilon_n$ &
$\vtwo\binom{2n}{n}$ & $\sum_j\lfloor n/2^j\rfloor$ & parity\\
\midrule
0  & 0000 & 0 & 0 &  1 & 0 &  0 & even\\
1  & 0001 & 1 & 1 & -1 & 1 &  1 & odd\\
2  & 0010 & 1 & 1 & -1 & 1 &  3 & odd\\
3  & 0011 & 2 & 0 &  1 & 2 &  4 & even\\
4  & 0100 & 1 & 1 & -1 & 1 &  7 & odd\\
5  & 0101 & 2 & 0 &  1 & 2 &  8 & even\\
6  & 0110 & 2 & 0 &  1 & 2 & 10 & even\\
7  & 0111 & 3 & 1 & -1 & 3 & 11 & odd\\
8  & 1000 & 1 & 1 & -1 & 1 & 15 & odd\\
9  & 1001 & 2 & 0 &  1 & 2 & 16 & even\\
10 & 1010 & 2 & 0 &  1 & 2 & 18 & even\\
11 & 1011 & 3 & 1 & -1 & 3 & 19 & odd\\
12 & 1100 & 2 & 0 &  1 & 2 & 22 & even\\
13 & 1101 & 3 & 1 & -1 & 3 & 23 & odd\\
14 & 1110 & 3 & 1 & -1 & 3 & 25 & odd\\
15 & 1111 & 4 & 0 &  1 & 4 & 26 & even\\
\bottomrule
\end{tabular}
\caption{The first sixteen terms and two arithmetic encodings developed below.}
\label{tab:first-values}
\end{table}

\begin{boxedremark}
\textbf{Proof-status convention.}
The source development already contains the binary recurrences, the finite
block product, Prouhet annihilation, the first surviving moment, iterated prefix
sums, and the binomial--logarithm formula.  They are recalled only when needed as
inputs or comparison points.  The valuation, Boolean-Möbius, determinant,
finite-Fourier, and Dirichlet-transform formulae below are proved mathematically
here, but this report does not silently promote them to theorem names in Lean.
\end{boxedremark}

\section{Binary recurrences, blocks, and uniqueness}
\label{sec:binary-recursions}

The two one-bit recurrences are the smallest complete description of the
sequence.

\begin{theorem}[Binary recursion and uniqueness]
\label{thm:binary-recursion}
The signs satisfy
\begin{equation}
  \varepsilon_0=1,
  \qquad
  \varepsilon_{2n}=\varepsilon_n,
  \qquad
  \varepsilon_{2n+1}=-\varepsilon_n.                       \label{eq:sign-recursion}
\end{equation}
Equivalently,
\begin{equation}
  \tau(0)=0,
  \qquad
  \tau(2n)=\tau(n),
  \qquad
  \tau(2n+1)=1-\tau(n).                                    \label{eq:bit-recursion}
\end{equation}
Either system, together with its initial value, determines its sequence uniquely.
\end{theorem}

\begin{proof}
Appending a terminal binary digit $0$ does not change the digit sum, while
appending a terminal $1$ increases it by one:
\[
  \wt(2n)=\wt(n),\qquad \wt(2n+1)=\wt(n)+1.
\]
This gives both recurrences.  For uniqueness, every positive integer is either
$2q$ or $2q+1$ with $q<n$, so strong induction reduces its value to a smaller
index and eventually to $0$.
\end{proof}

The recurrences can be iterated by more than one binary digit at a time.

\begin{theorem}[Binary block concatenation]
\label{thm:block-concatenation}
For $m,q\ge0$ and $0\le r<2^m$,
\begin{align}
  \wt(2^m q+r)&=\wt(q)+\wt(r),                              \label{eq:weight-block}\\
  \varepsilon_{2^m q+r}&=\varepsilon_q\varepsilon_r,       \label{eq:sign-block}\\
  \tau(2^m q+r)&=\tau(q)\xor\tau(r).                       \label{eq:tau-block}
\end{align}
Here the last $\xor$ is ordinary XOR of two bits.
\end{theorem}

\begin{proof}
Because $r<2^m$, multiplication of $q$ by $2^m$ shifts its binary expansion by
$m$ places and leaves those $m$ places zero.  Adding $r$ therefore produces no
carry and simply concatenates the two binary blocks.  The weight is additive;
the sign and parity statements follow.
\end{proof}

Let
\begin{equation}
  W_m=\tau(0)\tau(1)\cdots\tau(2^m-1)                       \label{eq:word-Wm}
\end{equation}
be the level-$m$ binary word, and write $\overline{W}$ for bitwise complement.
The case $q=0,1$ of \cref{thm:block-concatenation} gives
\begin{equation}
  W_{m+1}=W_m\,\overline{W_m}.                              \label{eq:word-substitution}
\end{equation}
Thus the infinite word is the fixed point beginning in $0$ of the length-two
morphism
\begin{equation}
  0\longmapsto01,
  \qquad
  1\longmapsto10.                                          \label{eq:thue-morphism}
\end{equation}
This classical description is not a separate coincidence: it is exactly the
block-concatenation law written at the level of words.

Two useful symmetries are immediate.

\begin{corollary}[Complement and bit-permutation symmetry]
\label{cor:complement-symmetry}
For $0\le n<2^m$,
\begin{equation}
  \varepsilon_{2^m-1-n}=(-1)^m\varepsilon_n,
  \qquad
  \tau(2^m-1-n)\equiv m+\tau(n)\pmod2.                     \label{eq:block-complement}
\end{equation}
Moreover, permuting the first $m$ binary digits of $n$ leaves $\varepsilon_n$
and $\tau(n)$ unchanged.
\end{corollary}

\begin{proof}
Within $m$ bits, $2^m-1-n$ is the bitwise complement of $n$, so its weight is
$m-\wt(n)$.  A permutation of digit positions leaves their sum unchanged.
\end{proof}

\section{Digit products, floor sums, and trigonometric encodings}
\label{sec:digit-floor}

\subsection{A product over the binary digits}

Since each digit $b\in\{0,1\}$ satisfies $(-1)^b=1-2b$, the definition itself
factorizes.

\begin{proposition}[Digit-product formula]
\label{prop:digit-product}
For every $n\ge0$,
\begin{equation}
  \boxed{\displaystyle
  \varepsilon_n=\prod_{j\ge0}\bigl(1-2b_j(n)\bigr),
  \qquad
  \tau(n)=\frac{1-\prod_{j\ge0}(1-2b_j(n))}{2}.}            \label{eq:digit-product}
\end{equation}
Only finitely many factors differ from $1$.
\end{proposition}

\begin{proof}
Multiply $(-1)^{b_j(n)}=1-2b_j(n)$ over all digit positions and use
$\sum_jb_j(n)=\wt(n)$.
\end{proof}

The individual digit can itself be written without a remainder operation.

\begin{lemma}[Floor extraction of one bit]
\label{lem:floor-bit}
For every $j,n\ge0$,
\begin{equation}
  b_j(n)=\left\lfloor\frac{n}{2^j}\right\rfloor
          -2\left\lfloor\frac{n}{2^{j+1}}\right\rfloor.     \label{eq:floor-bit}
\end{equation}
Consequently
\begin{equation}
  (-1)^{b_j(n)}=(-1)^{\lfloor n/2^j\rfloor}.               \label{eq:floor-bit-sign}
\end{equation}
\end{lemma}

\begin{proof}
Write $\lfloor n/2^j\rfloor=2q+r$ with $r\in\{0,1\}$.  Then
$q=\lfloor n/2^{j+1}\rfloor$ and $r=b_j(n)$.  The sign identity follows because
the two integers differ by an even number.
\end{proof}

\subsection{Legendre's floor sum before factorials}

The following elementary identity will reappear as the $p=2$ case of Legendre's
factorial formula.

\begin{lemma}[Dyadic floor sum]
\label{lem:dyadic-floor-sum}
For every $n\ge0$,
\begin{equation}
  \sum_{j\ge1}\left\lfloor\frac{n}{2^j}\right\rfloor
   =n-\wt(n),                                               \label{eq:floor-sum-legendre}
\end{equation}
so that
\begin{equation}
  \sum_{j\ge0}\left\lfloor\frac{n}{2^j}\right\rfloor
   =2n-\wt(n).                                              \label{eq:floor-sum-total}
\end{equation}
\end{lemma}

\begin{proof}
Insert \eqref{eq:binary-expansion} and interchange finite sums:
\begin{align*}
 \sum_{j\ge1}\left\lfloor\frac n{2^j}\right\rfloor
 &=\sum_{j\ge1}\sum_{\ell\ge j}b_\ell(n)2^{\ell-j}
  =\sum_{\ell\ge1}b_\ell(n)\sum_{j=1}^{\ell}2^{\ell-j}\\
 &=\sum_{\ell\ge1}b_\ell(n)(2^\ell-1)
  =n-\sum_{\ell\ge0}b_\ell(n)=n-\wt(n).
\end{align*}
Adding the $j=0$ term $n$ gives the second identity.
\end{proof}

Taking parity in \eqref{eq:floor-sum-total} eliminates $2n$.

\begin{theorem}[Floor-sum formulae]
\label{thm:floor-formulas}
For every $n\ge0$,
\begin{equation}
  \boxed{\displaystyle
  \varepsilon_n=
  (-1)^{\sum_{j\ge0}\lfloor n/2^j\rfloor}}
  \qquad\text{and}\qquad
  \tau(n)\equiv
  \sum_{j\ge0}\left\lfloor\frac{n}{2^j}\right\rfloor
  \pmod2.}                                                  \label{eq:floor-formulas}
\end{equation}
Equivalently,
\begin{equation}
  \varepsilon_n=\prod_{j\ge0}(-1)^{\lfloor n/2^j\rfloor}. \label{eq:floor-product}
\end{equation}
\end{theorem}

\begin{proof}
Equation~\eqref{eq:floor-sum-total} says that the exponent on the right differs
from $\wt(n)$ by the even integer $2n-2\wt(n)$.  The product form is simply
$(-1)^{a+b}=(-1)^a(-1)^b$.
\end{proof}

Formula~\eqref{eq:floor-formulas} uses no digit extraction and no recursion.  If
$n>0$, the sum may be stopped at $j=\lfloor\log_2n\rfloor$; for $n=0$ every term
is zero.

\subsection{A finite trigonometric sign product}

The square-wave identity
\begin{equation}
  \sgn(\sin\pi x)=(-1)^{\lfloor x\rfloor}
  \qquad(x\notin\Z)                                        \label{eq:sine-square-wave}
\end{equation}
turns the floor product into an analytic-looking formula.

\begin{theorem}[Trigonometric sign formula]
\label{thm:trig-sign}
Let $m\ge1$ and $0\le n<2^m$.  Then no factor below vanishes and
\begin{equation}
  \boxed{\displaystyle
  \varepsilon_n=
  \prod_{j=0}^{m-1}
  \sgn\!\left(
     \sin\frac{\pi(n+\tfrac12)}{2^j}
  \right).}                                                 \label{eq:trig-sign}
\end{equation}
In fact the $j$th factor is exactly $(-1)^{b_j(n)}$, and hence
\begin{equation}
  b_j(n)=\frac{1-
  \sgn\!\left(\sin\bigl(\pi(n+\tfrac12)/2^j\bigr)\right)}2
  \qquad(0\le j<m).                                        \label{eq:trig-bit}
\end{equation}
\end{theorem}

\begin{proof}
Because $n+\tfrac12$ is a half-integer, $(n+\tfrac12)/2^j$ is never an integer.
Moreover
\[
 \left\lfloor\frac{n+\tfrac12}{2^j}\right\rfloor
 =\left\lfloor\frac n{2^j}\right\rfloor.
\]
Apply \eqref{eq:sine-square-wave}, then
\eqref{eq:floor-bit-sign}; multiplying over $0\le j<m$ gives
\eqref{eq:trig-sign} because all higher digits vanish.
\end{proof}

\begin{remark}
The appearance of trigonometric functions is purely an encoding of dyadic
square waves; it does not smooth the sequence.  The half-unit shift is important:
without it, some sine factors would vanish at integer arguments and the sign
would be undefined.
\end{remark}

\section{Factorial valuations, central binomial coefficients, and carries}
\label{sec:valuations}

\subsection{A one-binomial-coefficient formula}

The floor sum in \cref{lem:dyadic-floor-sum} is exactly Legendre's formula for
the exponent of $2$ in a factorial.

\begin{theorem}[Factorial-valuation formula]
\label{thm:factorial-valuation}
For every $n\ge0$,
\begin{equation}
  \vtwo(n!)=\sum_{j\ge1}\left\lfloor\frac n{2^j}\right\rfloor
            =n-\wt(n).                                     \label{eq:legendre-two}
\end{equation}
Consequently
\begin{equation}
  \boxed{\displaystyle
  \varepsilon_n=(-1)^{n-\vtwo(n!)}=(-1)^{n+\vtwo(n!)},
  \qquad
  \tau(n)\equiv n+\vtwo(n!)\pmod2.}                       \label{eq:factorial-sign}
\end{equation}
\end{theorem}

\begin{proof}
Every multiple of $2$ contributes at least one factor of $2$ to $n!$, every
multiple of $4$ contributes one additional factor, and so on.  This gives the
floor sum, which equals $n-\wt(n)$ by \cref{lem:dyadic-floor-sum}.  The two
exponents $n-\vtwo(n!)$ and $n+\vtwo(n!)$ have the same parity.
\end{proof}

The central binomial coefficient removes the extra visible $n$ altogether.

\begin{theorem}[Central-binomial valuation]
\label{thm:central-binomial}
For every $n\ge0$,
\begin{equation}
  \boxed{\displaystyle
  \vtwo\binom{2n}{n}=\wt(n).}                               \label{eq:central-v2-weight}
\end{equation}
Hence
\begin{equation}
  \boxed{\displaystyle
  \varepsilon_n=(-1)^{\vtwo\binom{2n}{n}},
  \qquad
  \tau(n)=\vtwo\binom{2n}{n}\bmod2.}                      \label{eq:central-thue-morse}
\end{equation}
\end{theorem}

\begin{proof}
Using \eqref{eq:legendre-two} and $\wt(2n)=\wt(n)$,
\begin{align*}
 \vtwo\binom{2n}{n}
 &=\vtwo((2n)!)-2\vtwo(n!)\\
 &=\bigl(2n-\wt(2n)\bigr)-2\bigl(n-\wt(n)\bigr)
  =\wt(n).
\end{align*}
Taking parity gives \eqref{eq:central-thue-morse}.
\end{proof}

\begin{boxedremark}
\textbf{The shortest Pascal-triangle replacement.}
The source formula \eqref{eq:source-binomial-log} inspects the parity of all
$n+1$ entries in row $n$ and then takes an exact logarithm.  Formula
\eqref{eq:central-thue-morse} instead inspects the two-adic order of the single
entry $\binom{2n}{n}$ in row $2n$.  It recovers not merely the parity but the
entire digit sum $\wt(n)$.
\end{boxedremark}

\subsection{Kummer's carry interpretation}

Kummer's theorem says that $\vtwo\binom{a+b}{a}$ is the number of carry
operations in the base-two addition of $a$ and $b$.  In the special case $a=b=n$,
\cref{thm:central-binomial} gives an exact and slightly surprising statement.

\begin{corollary}[Doubling carry count]
\label{cor:doubling-carries}
The total number of binary carry operations performed when adding $n+n$ is
exactly $\wt(n)$.  Therefore $\tau(n)$ is the parity of that carry count.
\end{corollary}

The general addition law is itself a Thue--Morse formula.

\begin{theorem}[Addition and carry cocycle]
\label{thm:addition-carry}
For all $a,b\ge0$,
\begin{equation}
  \vtwo\binom{a+b}{a}
  =\wt(a)+\wt(b)-\wt(a+b),                                 \label{eq:kummer-digit-sum}
\end{equation}
and hence
\begin{equation}
  \boxed{\displaystyle
  \varepsilon_{a+b}
   =\varepsilon_a\varepsilon_b
      (-1)^{\vtwo\binom{a+b}{a}}.}                         \label{eq:addition-cocycle}
\end{equation}
Equivalently,
\begin{equation}
  \tau(a+b)\equiv\tau(a)+\tau(b)+
  \vtwo\binom{a+b}{a}\pmod2.                              \label{eq:addition-bit}
\end{equation}
\end{theorem}

\begin{proof}
Legendre's formula gives
\begin{align*}
 \vtwo\binom{a+b}{a}
 &=(a+b-\wt(a+b))-(a-\wt(a))-(b-\wt(b))\\
 &=\wt(a)+\wt(b)-\wt(a+b).
\end{align*}
Rearrange, exponentiate $-1$, and use
$(-1)^{-r}=(-1)^r$ for integral $r$.
\end{proof}

Thus ordinary addition differs from carry-free XOR by the sign of a binomial
valuation.  The same argument gives a multinomial extension.

\begin{corollary}[Multinomial carry formula]
\label{cor:multinomial-carry}
For $a_1,\ldots,a_r\ge0$ and $N=a_1+\cdots+a_r$,
\begin{equation}
  \varepsilon_N
  =\left(\prod_{j=1}^{r}\varepsilon_{a_j}\right)
   (-1)^{\vtwo\binom{N}{a_1,\ldots,a_r}},                  \label{eq:multinomial-sign}
\end{equation}
where
\begin{equation}
 \vtwo\binom{N}{a_1,\ldots,a_r}
 =\sum_{j=1}^{r}\wt(a_j)-\wt(N).                           \label{eq:multinomial-v2}
\end{equation}
\end{corollary}

\subsection{A local formula from the ruler function}

Incrementing $n$ flips all trailing $1$-bits to $0$ and then flips the next
$0$-bit to $1$.  The number of trailing $1$-bits of $n$ is $\vtwo(n+1)$.

\begin{theorem}[Adjacent-term valuation formula]
\label{thm:adjacent-term}
For every $n\ge0$,
\begin{align}
  \wt(n+1)-\wt(n)&=1-\vtwo(n+1),                            \label{eq:weight-increment}\\
  \varepsilon_{n+1}&=(-1)^{1+\vtwo(n+1)}\varepsilon_n,     \label{eq:sign-increment}\\
  (-1)^{\vtwo(n+1)}&=-\varepsilon_n\varepsilon_{n+1}.      \label{eq:ruler-from-adjacent}
\end{align}
Consequently
\begin{equation}
  \boxed{\displaystyle
  \varepsilon_n=
  \prod_{k=1}^{n}(-1)^{1+\vtwo(k)}
  =(-1)^{n+\sum_{k=1}^{n}\vtwo(k)}.}                       \label{eq:ruler-product}
\end{equation}
\end{theorem}

\begin{proof}
Let $r=\vtwo(n+1)$.  Then the binary expansion of $n$ ends in exactly $r$ ones.
The transition $n\mapsto n+1$ deletes those $r$ ones and inserts one new one,
proving \eqref{eq:weight-increment}.  Taking $(-1)$ to both sides gives
\eqref{eq:sign-increment}; multiplying by $-\varepsilon_n$ gives
\eqref{eq:ruler-from-adjacent}.  Finally telescope the adjacent ratios from
$\varepsilon_0=1$ to $\varepsilon_n$.
\end{proof}

Since $\sum_{k=1}^{n}\vtwo(k)=\vtwo(n!)$, the product formula closes the circle
back to \eqref{eq:factorial-sign}.  It also identifies the parity of the classical
ruler sequence as the multiplicative first difference of Thue--Morse.

\section{XOR characters, Lucas parity, and Boolean Möbius inversion}
\label{sec:boolean}

\subsection{The parity character of the Boolean group}

On nonnegative integers equipped with bitwise XOR, binary weight is additive
modulo two.

\begin{theorem}[XOR character]
\label{thm:xor-character}
For all $a,b\ge0$,
\begin{align}
  \tau(a\xor b)&=\tau(a)\xor\tau(b),                        \label{eq:tau-xor}\\
  \boxed{\displaystyle
  \varepsilon_{a\xor b}=\varepsilon_a\varepsilon_b.}      \label{eq:epsilon-xor}
\end{align}
Thus $n\mapsto\varepsilon_n$ is a group character from the direct sum of
countably many copies of $\Z/2\Z$ to $\{\pm1\}$.
\end{theorem}

\begin{proof}
At every digit position, XOR adds the two bits modulo two.  Summing those digit
equalities modulo two gives
$\wt(a\xor b)\equiv\wt(a)+\wt(b)\pmod2$, and both formulae follow.
\end{proof}

The block law \eqref{eq:sign-block} is the no-overlap special case of
\eqref{eq:epsilon-xor}, because $2^mq+r=(2^mq)\xor r$ for $r<2^m$.
Several finite-cube identities now become one-line consequences.

\begin{corollary}[XOR autocorrelation and Hamming spheres]
\label{cor:xor-correlations}
Fix $m\ge0$ and $0\le a<2^m$.  Then
\begin{equation}
  \sum_{n=0}^{2^m-1}\varepsilon_n\varepsilon_{n\xor a}
  =2^m\varepsilon_a.                                       \label{eq:xor-autocorrelation}
\end{equation}
More generally, for $0\le r\le m$ and $0\le a<2^m$,
\begin{equation}
  \sum_{\substack{0\le n<2^m\\ \wt(n\xor a)=r}}
       \varepsilon_n
  =\varepsilon_a(-1)^r\binom mr.                           \label{eq:hamming-sphere-sum}
\end{equation}
\end{corollary}

\begin{proof}
Equation~\eqref{eq:epsilon-xor} gives
$\varepsilon_n\varepsilon_{n\xor a}=\varepsilon_a$ termwise.
For the second formula put $h=n\xor a$.  Then
$\varepsilon_n=\varepsilon_a\varepsilon_h=\varepsilon_a(-1)^r$,
and there are $\binom mr$ masks $h$ of weight $r$.
\end{proof}

\subsection{Submasks and Pascal parity}

The relation $k\preceqbits n$ identifies the set bits of an integer with a finite
Boolean lattice.  Lucas's theorem modulo two says precisely that this lattice is
visible in Pascal's triangle.

\begin{theorem}[Lucas's submask criterion]
\label{thm:lucas-submask}
For $0\le k\le n$,
\begin{equation}
  \binom nk\equiv1\pmod2
  \quad\Longleftrightarrow\quad
  k\preceqbits n.                                          \label{eq:lucas-submask}
\end{equation}
Equivalently,
\begin{equation}
  \ind_{k\preceqbits n}
  =\frac{1-(-1)^{\binom nk}}{2}.                            \label{eq:odd-indicator}
\end{equation}
\end{theorem}

\begin{proof}
Write $n=\sum n_j2^j$ and $k=\sum k_j2^j$.  Lucas's congruence gives
\[
  \binom nk\equiv\prod_{j\ge0}\binom{n_j}{k_j}\pmod2.
\]
A factor is zero exactly when $n_j=0$ and $k_j=1$; otherwise it is one.  Hence
the product is one exactly when every $1$-bit of $k$ occurs at a $1$-bit of $n$.
The indicator formula distinguishes odd and even integers by the sign
$(-1)^m$.
\end{proof}

The following polynomial identity simultaneously contains Glaisher's count,
Boolean cancellation, and a new reconstruction formula.

\begin{theorem}[Submask weight enumerator]
\label{thm:submask-enumerator}
For every $n\ge0$ and every indeterminate $z$,
\begin{equation}
  \boxed{\displaystyle
  \sum_{k\preceqbits n}z^{\wt(k)}=(1+z)^{\wt(n)}.}          \label{eq:submask-enumerator}
\end{equation}
Equivalently,
\begin{equation}
  \boxed{\displaystyle
  z^{\wt(n)}=\sum_{k\preceqbits n}(z-1)^{\wt(k)}.}         \label{eq:boolean-binomial-inversion}
\end{equation}
\end{theorem}

\begin{proof}
A submask $k$ is obtained by choosing an arbitrary subset of the $\wt(n)$
positions at which $n$ has a $1$.  Choosing a position contributes one factor
$z$, while omitting it contributes $1$, giving $(1+z)^{\wt(n)}$.  Replace $z$
by $z-1$ in the first identity to obtain the second.
\end{proof}

Important specializations are worth separating.

\begin{corollary}[Three submask formulae]
\label{cor:submask-specializations}
For every $n\ge0$,
\begin{align}
  \#\{k:k\preceqbits n\}&=2^{\wt(n)},                       \label{eq:glaisher-submask}\\
  \sum_{k\preceqbits n}\varepsilon_k
    &=\begin{cases}1,&n=0,\\0,&n>0,\end{cases}              \label{eq:boolean-mobius-cancel}\\
  \boxed{\displaystyle
  \varepsilon_n&=\sum_{k\preceqbits n}(-2)^{\wt(k)}.}      \label{eq:submask-reconstruct}
\end{align}
\end{corollary}

\begin{proof}
Set $z=1$ and $z=-1$ in \eqref{eq:submask-enumerator}; then set $z=-1$ in
\eqref{eq:boolean-binomial-inversion}.
\end{proof}

Formula~\eqref{eq:submask-reconstruct} is a Boolean-binomial transform rather
than an ordinary binomial sum.  Lucas's criterion and the central-binomial
valuation now convert it into a formula involving only ordinary binomial
coefficients and $\vtwo$.

\begin{theorem}[A Pascal-parity and central-valuation formula]
\label{thm:pure-binomial-formula}
For every $n\ge0$,
\begin{equation}
 \boxed{\displaystyle
 \varepsilon_n=
 \sum_{k=0}^{n}
   \frac{1-(-1)^{\binom nk}}{2}
   (-2)^{\vtwo\binom{2k}{k}}.}                             \label{eq:pure-binomial-formula}
\end{equation}
Consequently
\begin{equation}
 \tau(n)=\frac12\left(
 1-
 \sum_{k=0}^{n}
   \frac{1-(-1)^{\binom nk}}{2}
   (-2)^{\vtwo\binom{2k}{k}}
 \right).                                                   \label{eq:pure-binomial-tau}
\end{equation}
\end{theorem}

\begin{proof}
The first factor in each summand is $\ind_{k\preceqbits n}$ by
\eqref{eq:odd-indicator}; the exponent in the second factor is $\wt(k)$ by
\eqref{eq:central-v2-weight}.  Thus the sum is exactly
\eqref{eq:submask-reconstruct}.  Apply \eqref{eq:normalization-bridge} for the
zero--one form.
\end{proof}

Unlike the source formula, \eqref{eq:pure-binomial-formula} uses no logarithm.
Its price is that it mixes two Pascal rows: row $n$ determines which $k$ are
submasks, and the central entry of row $2k$ determines the weight attached to
that submask.

\subsection{The Möbius function of the bitwise divisibility poset}

The interval $[k,n]$ in the submask poset is a Boolean lattice on the bit
positions that occur in $n$ but not in $k$.  Its Möbius function is therefore
an appropriate Thue--Morse sign.

\begin{theorem}[Boolean Möbius formula]
\label{thm:boolean-mobius}
If $k\preceqbits n$, then the Möbius function of the submask poset is
\begin{equation}
  \mu_2(k,n)=(-1)^{\wt(n)-\wt(k)}
             =\varepsilon_{n\xor k}.                       \label{eq:boolean-mobius}
\end{equation}
Accordingly, for functions $f,g$ on the nonnegative integers,
\begin{equation}
  g(n)=\sum_{k\preceqbits n}f(k)                            \label{eq:boolean-zeta}
\end{equation}
if and only if
\begin{equation}
  \boxed{\displaystyle
  f(n)=\sum_{k\preceqbits n}
        \varepsilon_{n\xor k}\,g(k).}                      \label{eq:boolean-mobius-inversion}
\end{equation}
\end{theorem}

\begin{proof}
If $k\preceqbits n$, then $n\xor k$ consists exactly of the bits in $n$ that are
absent from $k$, so its weight is $\wt(n)-\wt(k)$.  The interval $[k,n]$ is the
Boolean lattice on those positions, whose Möbius function is $(-1)^r$ in rank
$r$.  Standard Möbius inversion gives the equivalence.
\end{proof}

This can be written as an inverse for the parity shadow of Pascal's triangle.
Let
\begin{equation}
  P_{n,k}=\ind_{\binom nk\text{ odd}}
  =\ind_{k\preceqbits n}\qquad(0\le k\le n).               \label{eq:pascal-parity-matrix}
\end{equation}
Then the locally finite lower-triangular inverse is
\begin{equation}
  \boxed{\displaystyle
  (P^{-1})_{n,k}=P_{n,k}(-1)^{\wt(n)-\wt(k)}
  =P_{n,k}\varepsilon_{n\xor k}.}                          \label{eq:pascal-parity-inverse}
\end{equation}
In particular, the first column of $P^{-1}$ is precisely the signed
Thue--Morse sequence:
\begin{equation}
  (P^{-1})_{n,0}=\varepsilon_n.                             \label{eq:pascal-inverse-first-column}
\end{equation}
This is a second, conceptually different way in which Pascal parity encodes
Thue--Morse: not by counting one row, but as the Möbius inverse of the entire
submask incidence matrix.

\subsection{Walsh--Hadamard localization}

Fix $m$ and identify $0\le n<2^m$ with the Boolean vector of its first $m$ bits.
For $r<2^m$ define the Walsh character
\begin{equation}
  \chi_r(n)=(-1)^{\wt(r\bitand n)}.                         \label{eq:walsh-character}
\end{equation}
The all-ones mask $\mathbf1_m=2^m-1$ selects every bit, so
\begin{equation}
  \varepsilon_n=\chi_{2^m-1}(n)\qquad(0\le n<2^m).         \label{eq:epsilon-walsh}
\end{equation}
Orthogonality of Boolean characters therefore gives a delta-function transform.

\begin{theorem}[Walsh transform of a Thue--Morse block]
\label{thm:walsh-transform}
For $0\le r<2^m$,
\begin{equation}
 \boxed{\displaystyle
 \sum_{n=0}^{2^m-1}
 (-1)^{\wt(n)+\wt(r\bitand n)}
 =
 \begin{cases}
   2^m,&r=2^m-1,\\
   0,&r\ne2^m-1.
 \end{cases}}                                               \label{eq:walsh-transform}
\end{equation}
\end{theorem}

\begin{proof}
The summand is $\chi_{2^m-1}(n)\chi_r(n)
=\chi_{(2^m-1)\xor r}(n)$.  A nontrivial Boolean character takes each sign
exactly $2^{m-1}$ times, while the trivial character is identically one.
\end{proof}

Thus, on the Boolean cube, Thue--Morse is not spectrally complicated at all: it
is a single Walsh frequency, namely the highest-order parity character.  The
complication seen by ordinary Fourier analysis comes from replacing XOR by
cyclic addition.

\section{Generating products, a valuation recurrence, Bell polynomials, and determinants}
\label{sec:generating}

\subsection{Finite and infinite lacunary products}

Define the level-$m$ block polynomial
\begin{equation}
  K_m(x)=\sum_{n=0}^{2^m-1}\varepsilon_nx^n.                \label{eq:Km-definition}
\end{equation}
The finite product below is already one of the central Thue--Morse identities in
the source development; we retain it because several new formulae follow from
its logarithm and its finite Fourier transform.

\begin{theorem}[Lacunary block product]
\label{thm:lacunary-block-product}
For every $m\ge0$,
\begin{equation}
  \boxed{\displaystyle
  K_m(x)=\prod_{j=0}^{m-1}\bigl(1-x^{2^j}\bigr).}           \label{eq:Km-product}
\end{equation}
In particular,
\begin{equation}
  K_{m+1}(x)=\bigl(1-x^{2^m}\bigr)K_m(x).                  \label{eq:Km-recursion}
\end{equation}
\end{theorem}

\begin{proof}
Expand the product.  Choosing the term $-x^{2^j}$ from the $j$th factor and $1$
from the other factors selects a subset $S\subseteq\{0,\ldots,m-1\}$ and
contributes
\[
  (-1)^{|S|}x^{\sum_{j\in S}2^j}.
\]
Every integer $0\le n<2^m$ has exactly one such subset, namely the support of its
binary expansion; the coefficient is $(-1)^{\wt(n)}=\varepsilon_n$.
\end{proof}

The product displays both reversal symmetry and all cyclotomic multiplicities.

\begin{proposition}[Reciprocity and cyclotomic factorization]
\label{prop:cyclotomic-factorization}
Let $N=2^m$ and $D=N-1$.  Then
\begin{equation}
  x^D K_m(x^{-1})=(-1)^mK_m(x),                             \label{eq:Km-reciprocity}
\end{equation}
and
\begin{equation}
  K_m(x)=(1-x)^m
  \prod_{q=0}^{m-2}\bigl(1+x^{2^q}\bigr)^{m-1-q}.          \label{eq:Km-cyclotomic}
\end{equation}
Equivalently, in terms of cyclotomic polynomials,
\begin{equation}
  K_m(x)=(-1)^m\prod_{q=0}^{m-1}
             \Phi_{2^q}(x)^{m-q},                          \label{eq:Km-Phi}
\end{equation}
where $\Phi_{2^0}=\Phi_1=x-1$.  A root of unity of exact order $2^q$ has
multiplicity $m-q$ in $K_m$ when $0\le q<m$, and is not a root when $q\ge m$.
\end{proposition}

\begin{proof}
For reciprocity,
\[
 x^D K_m(x^{-1})
 =\prod_{j=0}^{m-1}(x^{2^j}-1)=(-1)^mK_m(x).
\]
For the second factorization use
$1-x^{2^j}=(1-x)\prod_{q=0}^{j-1}(1+x^{2^q})$ and count how many indices
$j$ contain each factor.  The cyclotomic form follows from
$x^{2^j}-1=\prod_{q=0}^{j}\Phi_{2^q}(x)$.  Its exponents give the root
multiplicities.
\end{proof}

Coefficient comparison in \eqref{eq:Km-reciprocity} recovers
\eqref{eq:block-complement}.  The zero of order $m$ at $x=1$ is the generating
function shadow of Prouhet annihilation.

The stabilized ordinary generating function is
\begin{equation}
  K(x)=\sum_{n\ge0}\varepsilon_nx^n.                        \label{eq:K-infinite-definition}
\end{equation}
As a formal series, each coefficient stabilizes in $K_m$; analytically, the
product converges normally on compact subsets of $|x|<1$ because
$\sum_j|x|^{2^j}<\infty$.

\begin{theorem}[Mahler product and the zero--one generating function]
\label{thm:Mahler-product}
For $|x|<1$, and also coefficientwise as a formal series,
\begin{align}
  \boxed{\displaystyle
  K(x)&=\prod_{j\ge0}(1-x^{2^j}),}                          \label{eq:K-infinite-product}\\
  K(x)&=(1-x)K(x^2).                                        \label{eq:K-Mahler}
\end{align}
If
\begin{equation}
  T(x)=\sum_{n\ge0}\tau(n)x^n,                             \label{eq:T-definition}
\end{equation}
then
\begin{align}
  \boxed{\displaystyle
  T(x)&=\frac12\left(\frac1{1-x}-K(x)\right),}             \label{eq:T-from-K}\\
  T(x)&=(1-x)T(x^2)+\frac{x}{1-x^2}.                       \label{eq:T-Mahler}
\end{align}
\end{theorem}

\begin{proof}
The product is the coefficientwise limit of \eqref{eq:Km-product}; separating
the $j=0$ factor and reindexing gives \eqref{eq:K-Mahler}.  Summing
$\varepsilon_n=1-2\tau(n)$ gives \eqref{eq:T-from-K}.  Alternatively, split
$T$ into even and odd indices and use \eqref{eq:bit-recursion}:
\begin{align*}
 T(x)
 &=\sum_{n\ge0}\tau(n)x^{2n}
   +\sum_{n\ge0}(1-\tau(n))x^{2n+1}\\
 &=(1-x)T(x^2)+\frac{x}{1-x^2}.
\end{align*}
\end{proof}

For example, the Prouhet--Thue--Morse binary constant
\begin{equation}
  \mathfrak t=\sum_{n\ge0}\frac{\tau(n)}{2^{n+1}}          \label{eq:TM-constant-definition}
\end{equation}
has the product representation
\begin{equation}
  \boxed{\displaystyle
  \mathfrak t=\frac12-\frac14
    \prod_{j\ge0}\left(1-2^{-2^j}\right).}                \label{eq:TM-constant-product}
\end{equation}
This is simply one half of \eqref{eq:T-from-K} evaluated at $x=1/2$.

\subsection{The logarithm reveals the ruler sequence}

The product contains a second valuation sequence after taking a logarithm.  Put
\begin{equation}
  a_r=\sum_{\substack{j\ge0\\2^j\mid r}}2^j
     =1+2+\cdots+2^{\vtwo(r)}
     =2^{\vtwo(r)+1}-1
  \qquad(r\ge1).                                           \label{eq:ar-definition}
\end{equation}

\begin{theorem}[Logarithmic series and convolution recurrence]
\label{thm:logarithmic-recurrence}
For $|x|<1$,
\begin{equation}
  \boxed{\displaystyle
  \log K(x)=-\sum_{r\ge1}\frac{a_r}{r}x^r,
  \qquad a_r=2^{\vtwo(r)+1}-1.}                            \label{eq:log-K}
\end{equation}
Consequently, for every $n\ge1$,
\begin{equation}
  \boxed{\displaystyle
  n\varepsilon_n
   =-\sum_{r=1}^{n}\bigl(2^{\vtwo(r)+1}-1\bigr)
        \varepsilon_{n-r}.}                                \label{eq:valuation-convolution}
\end{equation}
Thus
\begin{equation}
  \varepsilon_n=-\frac1n\sum_{r=1}^{n}a_r\varepsilon_{n-r}
  \qquad(n\ge1),                                          \label{eq:valuation-recurrence}
\end{equation}
with $\varepsilon_0=1$.
\end{theorem}

\begin{proof}
Absolute convergence permits the rearrangement
\begin{align*}
 \log K(x)
 &=\sum_{j\ge0}\log(1-x^{2^j})
  =-\sum_{j\ge0}\sum_{q\ge1}\frac{x^{q2^j}}q.
\end{align*}
The coefficient of $x^r$ receives a contribution only when $2^j\mid r$; then
$q=r/2^j$ and the contribution is $-2^j/r$.  Summing over such $j$ gives
$-a_r/r$.  Differentiating,
\[
 \frac{xK'(x)}{K(x)}=-\sum_{r\ge1}a_rx^r.
\]
Multiply by $K(x)=\sum_{n\ge0}\varepsilon_nx^n$ and compare the coefficient of
$x^n$.
\end{proof}

Formula~\eqref{eq:valuation-recurrence} is quite unlike the binary recursion:
it computes the $n$th term from all earlier terms and the two-adic orders of the
ordinary indices $1,\ldots,n$.  The divisibility data enter only through the odd
integers $a_r=1,3,1,7,1,3,1,15,\ldots$.

\subsection{An explicit partition sum and Bell-polynomial formula}

The exponential of \eqref{eq:log-K} turns the recurrence into a finite closed
sum over integer partitions.

\begin{theorem}[Partition-sum formula]
\label{thm:partition-sum}
For every $n\ge0$,
\begin{equation}
 \boxed{\displaystyle
 \varepsilon_n=
 \sum_{\substack{k_1,\ldots,k_n\ge0\\
                  1k_1+2k_2+\cdots+nk_n=n}}
 \prod_{r=1}^{n}
 \frac1{k_r!}
 \left(-\frac{2^{\vtwo(r)+1}-1}{r}\right)^{k_r}.}          \label{eq:partition-sum}
\end{equation}
At $n=0$ the unique empty tuple contributes $1$.
\end{theorem}

\begin{proof}
From \eqref{eq:log-K},
\[
 K(x)=\prod_{r\ge1}
 \exp\left(-\frac{a_r}{r}x^r\right)
 =\prod_{r\ge1}\sum_{k_r\ge0}
   \frac1{k_r!}\left(-\frac{a_r}{r}\right)^{k_r}x^{rk_r}.
\]
To obtain the coefficient of $x^n$, only $1\le r\le n$ can occur, and the
indices must satisfy $\sum_r rk_r=n$.
\end{proof}

The summands in \eqref{eq:partition-sum} are generally rational, yet their sum
is always exactly $\pm1$.  This cancellation is not apparent termwise.

For a compact version, define the complete exponential Bell polynomials by
\begin{equation}
 \exp\left(\sum_{r\ge1}y_r\frac{x^r}{r!}\right)
 =\sum_{n\ge0}\Bell_n(y_1,\ldots,y_n)\frac{x^n}{n!}.       \label{eq:Bell-definition}
\end{equation}

\begin{corollary}[Bell-polynomial formula]
\label{cor:Bell-formula}
For every $n\ge0$,
\begin{equation}
 \boxed{\displaystyle
 \varepsilon_n=\frac1{n!}\Bell_n
 \bigl(-0!a_1,-1!a_2,\ldots,-(n-1)!a_n\bigr),
 \qquad a_r=2^{\vtwo(r)+1}-1.}                             \label{eq:Bell-formula}
\end{equation}
\end{corollary}

\begin{proof}
In \eqref{eq:Bell-definition} take $y_r=-(r-1)!a_r$, so that
$y_r/r!=-a_r/r$, and compare with \eqref{eq:log-K}.
\end{proof}

\subsection{A lower-Hessenberg determinant}

The same coefficient recurrence has a determinant solution.  For $n\ge1$, let
$H_n$ be the $n\times n$ lower-Hessenberg matrix
\begin{equation}
 H_n=\begin{pmatrix}
 a_1&1&0&0&\cdots&0\\
 a_2&a_1&2&0&\cdots&0\\
 a_3&a_2&a_1&3&\ddots&\vdots\\
 \vdots&\vdots&\ddots&\ddots&\ddots&0\\
 a_{n-1}&a_{n-2}&\cdots&a_2&a_1&n-1\\
 a_n&a_{n-1}&\cdots&a_3&a_2&a_1
 \end{pmatrix},                                            \label{eq:Hn-matrix}
\end{equation}
where $a_r=2^{\vtwo(r)+1}-1$.

\begin{theorem}[Determinant formula]
\label{thm:determinant-formula}
For every $n\ge1$,
\begin{equation}
  \boxed{\displaystyle
  \varepsilon_n=\frac{(-1)^n}{n!}\det H_n.}                \label{eq:determinant-formula}
\end{equation}
Equivalently, $\det H_n=(-1)^nn!\varepsilon_n$, so every one of these
nontrivial integer determinants has absolute value exactly $n!$.
\end{theorem}

\begin{proof}
Set $D_0=1$ and $D_n=\det H_n$.  Expansion of a lower-Hessenberg determinant
along the last row, or induction by continuant elimination, gives
\begin{equation}
 D_n=\sum_{r=1}^{n}(-1)^{r+1}
       \frac{(n-1)!}{(n-r)!}\,a_rD_{n-r}.                  \label{eq:Hessenberg-recurrence}
\end{equation}
Now put $c_n=(-1)^nD_n/n!$.  Multiplying
\eqref{eq:Hessenberg-recurrence} by $(-1)^n/n!$ yields
\[
 nc_n=-\sum_{r=1}^{n}a_rc_{n-r},\qquad c_0=1.
\]
This is exactly recurrence~\eqref{eq:valuation-convolution}, whose solution is
unique by induction.  Therefore $c_n=\varepsilon_n$.
\end{proof}

For example,
\begin{align*}
 \varepsilon_4
 &=\frac1{4!}
 \det\begin{pmatrix}
 1&1&0&0\\
 3&1&2&0\\
 1&3&1&3\\
 7&1&3&1
 \end{pmatrix}
 =\frac{-24}{24}=-1,                                      \label{eq:det-example}
\end{align*}
where $(a_1,a_2,a_3,a_4)=(1,3,1,7)$.  The determinant formula conceals all
binary digits inside the ruler values on its subdiagonals.


\section{Exact prefix sums, interval discrepancy, and dyadic blocks}
\label{sec:prefix-balance}

The signed sequence has exceptionally small partial sums.  Define the half-open
prefix sum
\begin{equation}
  S(N)=\sum_{0\le n<N}\varepsilon_n
  \qquad(N\ge0).                                           \label{eq:prefix-S-definition}
\end{equation}
Pairing consecutive terms evaluates it completely.

\begin{theorem}[Exact prefix-sum formula]
\label{thm:exact-prefix-sum}
For every $q\ge0$,
\begin{equation}
  \boxed{\displaystyle
  S(2q)=0,
  \qquad
  S(2q+1)=\varepsilon_q.}                                  \label{eq:exact-prefix-sum}
\end{equation}
Equivalently,
\begin{equation}
  S(N)=\ind_{N\text{ odd}}\,
       \varepsilon_{(N-1)/2}.                              \label{eq:prefix-indicator-form}
\end{equation}
In particular, $|S(N)|\le1$ for every $N$.
\end{theorem}

\begin{proof}
The binary recursion gives
$\varepsilon_{2j}+\varepsilon_{2j+1}=0$.  Hence the first $2q$ terms
cancel in $q$ pairs.  Appending the term at $2q$ leaves
$S(2q+1)=\varepsilon_{2q}=\varepsilon_q$.
\end{proof}

Let
\begin{equation}
  A(N)=\sum_{0\le n<N}\tau(n)                              \label{eq:prefix-A-definition}
\end{equation}
be the number of $1$-bits among the first $N$ values of the zero--one sequence.
Since $\tau=(1-\varepsilon)/2$, the preceding theorem gives an exact count.

\begin{corollary}[Exact zero--one balance]
\label{cor:exact-zero-one-balance}
For every $q\ge0$,
\begin{equation}
  \boxed{\displaystyle
  A(2q)=q,
  \qquad
  A(2q+1)=q+\tau(q).}                                      \label{eq:exact-one-count}
\end{equation}
Thus every even prefix contains exactly as many zeros as ones, and every odd
prefix differs from perfect balance by one symbol.
\end{corollary}

\begin{proof}
Use $A(N)=(N-S(N))/2$ and
$\tau(q)=(1-\varepsilon_q)/2$.
\end{proof}

The same formula evaluates every finite interval, not only prefixes.

\begin{corollary}[Closed formula for interval sums]
\label{cor:interval-sum}
For integers $0\le A\le B$,
\begin{equation}
 \boxed{\displaystyle
 \sum_{n=A}^{B-1}\varepsilon_n
 =\ind_{B\text{ odd}}\varepsilon_{(B-1)/2}
  -\ind_{A\text{ odd}}\varepsilon_{(A-1)/2}.}             \label{eq:interval-sum}
\end{equation}
Consequently
\begin{equation}
  \left|\sum_{n=A}^{B-1}\varepsilon_n\right|\le2,        \label{eq:interval-discrepancy-bound}
\end{equation}
and the number of $1$'s in any interval differs from half its length by at
most $1$.
\end{corollary}

\begin{proof}
The interval sum is $S(B)-S(A)$; substitute
\eqref{eq:prefix-indicator-form}.  The discrepancy statement follows from
$\sum\tau(n)=\bigl((B-A)-\sum\varepsilon_n\bigr)/2$.
\end{proof}

The bound $2$ is attained: for example,
$\varepsilon_1+\varepsilon_2=-2$.  For prefixes the stronger bound $1$ is
best possible.

\subsection{Self-similar partial blocks}

The concatenation law also gives every partial sum inside a dyadic block.

\begin{proposition}[Partial dyadic-block formula]
\label{prop:partial-dyadic-block}
Let $m,q\ge0$ and $0\le r\le2^m$.  Then
\begin{equation}
  \boxed{\displaystyle
  \sum_{j=0}^{r-1}\varepsilon_{2^m q+j}
  =\varepsilon_q S(r).}                                   \label{eq:partial-dyadic-block}
\end{equation}
In particular, when $m\ge1$,
\begin{equation}
  \sum_{j=0}^{2^m-1}\varepsilon_{2^m q+j}=0.              \label{eq:full-dyadic-block-zero}
\end{equation}
The corresponding number of $1$'s in the first $r$ positions of this block is
\begin{equation}
  \sum_{j=0}^{r-1}\tau(2^m q+j)
  =\frac{r-\varepsilon_qS(r)}2.                            \label{eq:partial-dyadic-one-count}
\end{equation}
\end{proposition}

\begin{proof}
For $j<2^m$, formula~\eqref{eq:sign-block} gives
$\varepsilon_{2^mq+j}=\varepsilon_q\varepsilon_j$.  Sum over
$0\le j<r$.  When $m\ge1$, $S(2^m)=0$.  The zero--one identity follows from
$\tau=(1-\varepsilon)/2$.
\end{proof}

Thus every dyadic block is either the basic level-$m$ sign word or its negative;
its internal discrepancy is the same function $S(r)$, up to one global sign.
This is a stronger statement than mere balance of complete blocks.

\subsection{The generating function of the prefix sums}

Cumulative summation has a particularly clean generating-series effect.

\begin{proposition}[Prefix generating series]
\label{prop:prefix-generating-series}
For $|x|<1$,
\begin{equation}
  \boxed{\displaystyle
  \sum_{N\ge0}S(N)x^N
   =\frac{x}{1-x}K(x)
   =xK(x^2)
   =x\prod_{j\ge1}(1-x^{2^j}).}                            \label{eq:prefix-generating-series}
\end{equation}
\end{proposition}

\begin{proof}
Absolute convergence for $|x|<1$ gives
\[
 \sum_{N\ge0}S(N)x^N
 =\sum_{n\ge0}\varepsilon_n\sum_{N\ge n+1}x^N
 =\frac{x}{1-x}K(x).
\]
The Mahler relation $K(x)=(1-x)K(x^2)$ gives the second expression.  Expanding
$xK(x^2)=\sum_{q\ge0}\varepsilon_qx^{2q+1}$ recovers
\eqref{eq:exact-prefix-sum} coefficient by coefficient.
\end{proof}

The last identity is the $k=1$ case of the iterated-prefix generating functions
already developed in \path{ThueMorseGenerating.lean}; it is included here
because, at one iteration, the coefficients collapse all the way back to a
subsequence of $\varepsilon$ itself.


\section{A master finite-difference formula}
\label{sec:finite-differences}

Many apparently unrelated identities follow from viewing the finite block
product as a polynomial in translation operators.  For $h\in\R$, let
\begin{equation}
  (E_hf)(x)=f(x+h),
  \qquad
  (\Delta_hf)(x)=f(x+h)-f(x)=(E_h-I)f(x).                  \label{eq:shift-difference}
\end{equation}
The operators commute because translations commute.

\begin{theorem}[Mixed finite-difference master identity]
\label{thm:finite-difference-master}
Let $m\ge0$ and let $f$ be any function defined at the integers
$0,1,\ldots,2^m-1$.  Then
\begin{equation}
 \boxed{\displaystyle
 \sum_{n=0}^{2^m-1}\varepsilon_nf(n)
 =\left[\prod_{j=0}^{m-1}(I-E_{2^j})f\right](0)
 =(-1)^m
   (\Delta_1\Delta_2\cdots\Delta_{2^{m-1}}f)(0).}         \label{eq:finite-difference-master}
\end{equation}
For $m=0$ both products are empty and the identity reads
$\varepsilon_0f(0)=f(0)$.
\end{theorem}

\begin{proof}
Expanding the product of shifts gives
\begin{align*}
 \left[\prod_{j=0}^{m-1}(I-E_{2^j})f\right](0)
 &=\sum_{J\subseteq\{0,\ldots,m-1\}}
   (-1)^{|J|}
   f\left(\sum_{j\in J}2^j\right).
\end{align*}
The subset $J$ is the support of the unique binary expansion of an integer
$n<2^m$, and $(-1)^{|J|}=\varepsilon_n$.  Finally,
$I-E_h=-\Delta_h$ for each of the $m$ factors.
\end{proof}

This is a useful ``formula generator'': substitute a class of test functions
whose mixed differences or derivatives can be evaluated.

\subsection{An exact integral representation}

Repeated use of the fundamental theorem of calculus converts every difference
into an integral.

\begin{theorem}[Cubical integral formula]
\label{thm:cubical-integral}
Let $m\ge1$, and suppose $f$ is $m$ times continuously differentiable on
$[0,2^m-1]$.  Then
\begin{equation}
 \boxed{\displaystyle
 \sum_{n=0}^{2^m-1}\varepsilon_nf(n)
 =(-1)^m
 \int_0^1\!\int_0^2\!\cdots\!\int_0^{2^{m-1}}
 f^{(m)}(u_0+u_1+\cdots+u_{m-1})
 \dd u_{m-1}\cdots\dd u_1\dd u_0.}                       \label{eq:cubical-integral}
\end{equation}
Consequently,
\begin{equation}
 \left|\sum_{n=0}^{2^m-1}\varepsilon_nf(n)\right|
 \le 2^{\binom m2}
     \max_{0\le x\le2^m-1}|f^{(m)}(x)|.                  \label{eq:cubical-integral-bound}
\end{equation}
\end{theorem}

\begin{proof}
For $h>0$,
\[
  (\Delta_hf)(x)=\int_0^h f'(x+u)\dd u.
\]
Iterate this identity with $h=1,2,\ldots,2^{m-1}$ and insert it into
\eqref{eq:finite-difference-master}.  The integration box has volume
\[
  \prod_{j=0}^{m-1}2^j=2^{0+1+\cdots+(m-1)}=2^{\binom m2},
\]
which gives the bound.
\end{proof}

After normalizing $u_j=2^jv_j$, formula~\eqref{eq:cubical-integral} can also be
written as
\begin{equation}
 \sum_{n<2^m}\varepsilon_nf(n)
 =(-1)^m2^{\binom m2}
 \int_{[0,1]^m}
 f^{(m)}\left(\sum_{j=0}^{m-1}2^jv_j\right)
 \dd v_0\cdots\dd v_{m-1}.                                \label{eq:unit-cube-integral}
\end{equation}
Thus the normalized signed sum is the expectation of $f^{(m)}$ at a finite
weighted sum of independent uniform variables.  This is the same dyadic
uniform-sum architecture that leads, after scaling and passage to a limit, to
the Fabius and Rvachev functions.

\subsection{Prouhet cancellation and the first surviving moment}

Taking $f$ to be a polynomial immediately gives the classical equal-sums-of-like-
powers property.

\begin{corollary}[Prouhet annihilation]
\label{cor:Prouhet-annihilation}
If $p$ is a polynomial of degree less than $m$, then
\begin{equation}
  \boxed{\displaystyle
  \sum_{n=0}^{2^m-1}\varepsilon_np(n)=0.}                  \label{eq:Prouhet-annihilation}
\end{equation}
Equivalently, for every integer $0\le d<m$,
\begin{equation}
 \sum_{\substack{0\le n<2^m\\\tau(n)=0}}n^d
 =
 \sum_{\substack{0\le n<2^m\\\tau(n)=1}}n^d.           \label{eq:Prouhet-two-classes}
\end{equation}
\end{corollary}

\begin{proof}
The $m$th derivative of $p$ is zero, so
\eqref{eq:cubical-integral} vanishes.  Splitting the signed sum according to
$\varepsilon_n=1$ and $-1$ gives the second form.
\end{proof}

At the first nonvanishing degree the integral is constant.

\begin{corollary}[First surviving power sum]
\label{cor:first-surviving-power-sum}
For every $m\ge0$,
\begin{equation}
 \boxed{\displaystyle
 \sum_{n=0}^{2^m-1}\varepsilon_nn^m
 =(-1)^m m!\,2^{\binom m2}.}                               \label{eq:first-surviving-power-sum}
\end{equation}
More generally, for every $a\in\R$,
\begin{equation}
 \sum_{n=0}^{2^m-1}\varepsilon_n(a+n)^m
 =(-1)^m m!\,2^{\binom m2};                               \label{eq:translated-first-moment}
\end{equation}
there is no dependence on the translation $a$.
\end{corollary}

\begin{proof}
For $f(x)=x^m$ or $f(x)=(a+x)^m$, one has $f^{(m)}=m!$.
Formula~\eqref{eq:cubical-integral} is therefore $(-1)^m m!$ times the volume
$2^{\binom m2}$.
\end{proof}

These are the polynomial identities already formalized in
\path{ThueMorsePrefix.lean} as
\lean{thueMorsePowerSum_eq_zero_of_lt} and
\lean{thueMorsePowerSum_self}.  The finite-difference theorem packages both
into one operator identity and extends them to arbitrary smooth test functions.

\subsection{Exponential, reciprocal-power, and logarithmic specializations}

For an exponential, each difference acts diagonally.

\begin{corollary}[Finite exponential transform]
\label{cor:finite-exponential-transform}
For every $z\in\C$ and $m\ge0$,
\begin{equation}
 \boxed{\displaystyle
 \sum_{n=0}^{2^m-1}\varepsilon_n\e^{zn}
 =\prod_{j=0}^{m-1}\left(1-\e^{2^jz}\right).}             \label{eq:finite-exponential-transform}
\end{equation}
\end{corollary}

\begin{proof}
Apply \eqref{eq:finite-difference-master} to $f(x)=\e^{zx}$, for which
$\Delta_hf=(\e^{zh}-1)f$.  The two signs cancel.  Equivalently, substitute
$x=\e^z$ in \eqref{eq:Km-product}.
\end{proof}

The reciprocal-power specialization has a less obvious and useful consequence:
the signed sum is always positive.

\begin{theorem}[Positive reciprocal-power formula]
\label{thm:positive-reciprocal-power}
Let $a>0$, $s>0$, and $m\ge1$.  With the rising factorial
$(s)_m=s(s+1)\cdots(s+m-1)$,
\begin{align}
 \boxed{\displaystyle
 \sum_{n=0}^{2^m-1}\frac{\varepsilon_n}{(a+n)^s}
 &=(s)_m
 \int_0^1\!\int_0^2\!\cdots\!\int_0^{2^{m-1}}
 \frac{\dd u_{m-1}\cdots\dd u_0}
 {(a+u_0+\cdots+u_{m-1})^{s+m}}}                           \label{eq:reciprocal-power-integral}\\
 &>0.                                                       \label{eq:reciprocal-power-positive}
 }
\end{align}
In particular, the even-parity indices in a complete dyadic block always
outweigh the odd-parity indices for every decreasing kernel $(a+n)^{-s}$.
\end{theorem}

\begin{proof}
For $f(x)=(a+x)^{-s}$,
\[
 f^{(m)}(x)=(-1)^m(s)_m(a+x)^{-s-m}.
\]
The factor $(-1)^m$ in \eqref{eq:cubical-integral} cancels this derivative
sign.  The resulting integrand is strictly positive.
\end{proof}

At $s=1$, repeated elementary integration also expresses the left-hand side as
a logarithm of a rational product; the integral form is usually cleaner and
makes positivity immediate.

Taking $f(x)=\log(a+x)$ gives such a product directly.

\begin{corollary}[A finite logarithmic product inequality]
\label{cor:finite-log-product}
Let $a>0$ and $m\ge1$.  Then
\begin{align}
 \sum_{n=0}^{2^m-1}\varepsilon_n\log(a+n)
 & =-(m-1)!
 \int_0^1\!\int_0^2\!\cdots\!\int_0^{2^{m-1}}
 \frac{\dd u_{m-1}\cdots\dd u_0}
 {(a+u_0+\cdots+u_{m-1})^m}                                \label{eq:finite-log-integral}\\
 &<0.                                                       \label{eq:finite-log-negative}
\end{align}
Consequently,
\begin{equation}
 \boxed{\displaystyle
 0<\prod_{n=0}^{2^m-1}(a+n)^{\varepsilon_n}<1.}            \label{eq:finite-log-product}
\end{equation}
\end{corollary}

\begin{proof}
For $m\ge1$,
\[
 \frac{\dd^m}{\dd x^m}\log(a+x)
 =(-1)^{m-1}\frac{(m-1)!}{(a+x)^m}.
\]
Multiplying by the outer $(-1)^m$ in
\eqref{eq:cubical-integral} gives the negative integral.  Exponentiation gives
\eqref{eq:finite-log-product}.
\end{proof}

Because $\sum_{n<2^m}\varepsilon_n=0$ for $m\ge1$, the product in
\eqref{eq:finite-log-product} is unchanged if every factor $a+n$ is multiplied
by the same positive constant.  The inequality is therefore genuinely about
the relative placement of the two Thue--Morse classes, not about units or
normalization.
